\begin{abstract}
 Robustness under distribution shift remains a central challenge in computer vision.
 While ensemble methods are widely used to improve generalization, their effectiveness depends on both model diversity and aggregation strategy.
 In this work, we compare homogeneous and heterogeneous ensembles under synthetic and real-world distribution shifts, and evaluate whether adaptive expert selection with a mixture-of-experts routing mechanism can better exploit model diversity

 We construct a homogeneous CNN ensemble and a heterogeneous ensemble combining convolutional networks with a Vision Transformer.
 Performance is evaluated on clean data, synthetic corruption benchmarks, and a student-collected out-of-distribution dataset.
 Oracle analysis is used to quantify the potential benefit of complementary expertise.

 Results show that homogeneous ensembles provide limited robustness gains due to correlated failures under shift.
 In contrast, heterogeneous ensembles exhibit larger oracle gaps and improved robustness in severe corruption regimes.
 However, a learned mixture-of-experts router primarily collapses to a dominant expert, yielding only modest improvements over uniform averaging.
 These findings indicate that ensemble robustness is constrained both by expert diversity and by the ability of routing mechanisms to effectively leverage that diversity.
\end{abstract}

